\subsection{Magnetic feedback that preserves the learned KAM dynamics}
\label{sec:full-fe-kam-control}
A spatial magnetic-field gradient supplies a tip-force input that, together
with the transverse field, preserves the certified conservative dynamics
despite physical damping. Both inputs act through the magnetic dipole
on the same finite-element beam.

Let $a_j$ be the physical mass-normalized modal positions of the complete
$64$-mode FE model, with velocities $v_j$, frequencies $\omega_j$,
tip-displacement participation $b_j$, and tip-rotation participation $t_j$.
The numerical theorem uses the exported modal coefficients and Hermite
shape map as the exact definition of the computational model. The
diagonalization residual relative to the assembled nodal matrices is
quantified below.
Write $\vartheta=\sum_jt_ja_j$ and $E_0=\mu B_0$. With transverse tip
force $F$ and transverse field $B_y$, the modal equations are
\begin{equation}
 \begin{aligned}
 \ddot a_j+\alpha\dot a_j+\omega_j^2a_j&=b_jF+t_jU,\\
 U&=E_0(\vartheta-\sin\vartheta)+\mu B_y\cos\vartheta.
 \end{aligned}
 \label{eq:full-fe-two-port}
\end{equation}
The frequencies already include the linear stiffness from $B_0$.
For the first two modes define
\begin{equation}
 A_{\rm mag}=\begin{pmatrix}b_1&t_1\\b_2&t_2\end{pmatrix}.
 \label{eq:magnetic-authority-matrix}
\end{equation}
An interval enclosure of its determinant excludes zero.

Let $h(X,Y)$ be the fitted Hamiltonian used by the KAM
certificate, including momentum-dependent quartics. Set
$a_i=S_iX_i$ and recover $Y$ from
\begin{equation}
 h_Y(X,Y)=\operatorname{diag}(S_i)^{-1}v_{1:2}/\omega_0.
 \label{eq:full-fe-legendre}
\end{equation}
This implicit velocity-to-momentum map cannot be replaced by $Y_i=v_i/(S_i
\omega_0\nu_i)$ when the learned momentum quartics are retained. The
verified lower bound on $h_{YY}$ below makes the map locally invertible and
injective in $Y$ on the stated real box. Define
\begin{align}
 a^{\rm des}_{i,tt}
 &=S_i\omega_0^2\bigl(h_{YX}h_Y-h_{YY}h_X\bigr)_i,\\
 \begin{pmatrix}F\\U\end{pmatrix}
 &=A_{\rm mag}^{-1}
   \left(a^{\rm des}_{1:2,tt}+\alpha v_{1:2}
          +\operatorname{diag}(\omega_1^2,\omega_2^2)a_{1:2}\right).
 \label{eq:full-fe-hamiltonian-feedback}
\end{align}
Consequently $\dot X=\omega_0h_Y$ and $\dot Y=-\omega_0h_X$ exactly
where the implicit inverse is defined. The proof follows by differentiating
\eqref{eq:full-fe-legendre}: the specified acceleration cancels the $h_{YX}$
term and invertibility of $h_{YY}$ gives the second canonical equation.

The physical commands realizing \eqref{eq:full-fe-hamiltonian-feedback} are
\begin{equation}
 \begin{aligned}
 B_y&=\frac{U-E_0(\vartheta-\sin\vartheta)}{\mu\cos\vartheta},\\
 G&=\frac{F}{\mu}
 \begin{pmatrix}-\sin\vartheta&\cos\vartheta\\
                 \cos\vartheta&\sin\vartheta\end{pmatrix}.
 \end{aligned}
 \label{eq:physical-gradient-command}
\end{equation}
Use $B(r)=B_{\rm tip}+G(r-r_{\rm tip})$ locally, with
$B_{\rm tip}=(B_0,B_y)$. The matrix $G$ is symmetric and traceless,
so this local planar field satisfies the source-free curl and divergence
constraints. For dipole $m=\mu(\cos\vartheta,\sin\vartheta)$ it gives
$Gm=(0,F)$, while the tip torque is
$\mu(B_y\cos\vartheta-B_0\sin\vartheta)$.
Moreover $\|G\|_2=|F|/\mu$. The applied mechanical power is
$F\dot w(L)+\mu B_y\cos\vartheta\,\dot\vartheta$, in addition to
the static magnetic energy already included in the Hamiltonian.
The construction assumes these local field and gradient commands are
available and the required modal and tip states are observed.

\subsection{Lifting a certified torus to the complete controlled FE model}
\label{sec:full-fe-kam-lift}
\begin{theorem}[Invariant-torus lift under physical magnetic feedback]
\label{thm:full-fe-kam-lift}
Suppose a real analytic embedded two-torus $K_\infty$ of $h$ has frequency
$\Omega\in\mathbb R^2$ in radians per physical second. Suppose $h_{YY}$
is positive definite on a neighbourhood of this torus with convex
$Y$-fibres for each $X$, $A_{\rm mag}$ is invertible,
$\alpha>0$, and the complete lifted solution has
$|\vartheta|<\pi/2$. Then the full FE closed loop
\eqref{eq:full-fe-two-port}--\eqref{eq:physical-gradient-command} has an
embedded invariant two-torus projecting to $K_\infty$. For a fixed low-mode
orbit on this torus, each high-mode solution converges exponentially to a
unique quasiperiodic response, provided the feedback remains in its
stated field and state domain.
\end{theorem}

\begin{proof}
The first two canonical modes follow $h$ exactly. Hence $F$ and $U$ in
\eqref{eq:full-fe-hamiltonian-feedback} depend only on their torus phase,
even though the physical field command also uses the complete tip angle.
For every $j\geq3$, the remaining dynamics are the stable linear equation
\eqref{eq:full-fe-two-port}, forced by analytic quasiperiodic functions.
If $F_k,U_k$ denote their Fourier coefficients in the convention
$e^{ik\cdot\phi}$, the unique response graph is
\begin{equation}
 a^*_{j,k}=
 \frac{b_jF_k+t_jU_k}
      {\omega_j^2-(k\cdot\Omega)^2+i\alpha(k\cdot\Omega)}.
 \label{eq:full-fe-lift-fourier}
\end{equation}
For real $k\cdot\Omega$, this denominator never vanishes: at zero
frequency its real part is $\omega_j^2>0$, and at nonzero frequency its
imaginary part is nonzero. Its inverse is uniformly bounded over real
frequencies. Exponential Fourier decay is therefore preserved on a
possibly smaller analytic strip. The velocity graph is obtained by
multiplication by $ik\cdot\Omega$ and is analytic on a smaller strip.
Positive definiteness on each convex $Y$-fibre makes
$Y\mapsto h_Y(X,Y)$ strictly monotone and injective. Thus the physical
low-mode positions and velocities preserve the injective canonical
embedding. Differences in the high modes obey the
homogeneous damped equations and decay exponentially.
\end{proof}

The lifted closed loop combines Hamiltonian low-mode motion with
contracting high-mode fibres. The contraction leaves each initial low-mode
torus invariant; it does not select a unique torus or make the full
controlled system Hamiltonian. Each lifted two-torus has zero volume in
the full $128$-dimensional phase space. The KAM-family measure therefore
refers to the conservative low-mode coordinates.

\paragraph{All-phase response and startup bounds.}
We bound both the response graph and the startup
with all high modal positions and velocities initially zero. Write
$h=h_2+R$, and let the certified torus satisfy
$\|K_\infty-K_0\|_\infty\leq\epsilon_K$, where
$K_0=0.1(\cos\phi_1,\cos\phi_2,-\sin\phi_1,-\sin\phi_2)$.
The first two components of the force demand can be written as
\begin{align}
 r_i={}&S_i\bigl\{\omega_0^2[
 (R_{YX}h_Y)_i-\nu_iR_{X_i}-(R_{YY}h_X)_i]\nonumber\\
 &\quad+\alpha\omega_0h_{Y_i}
 +(\omega_i^2-\omega_0^2\nu_i^2)X_i\bigr\},
 \label{eq:full-fe-small-force}
\end{align}
This identity cancels the large linear restoring terms before bounding.
The linear part on $K_0$ consists of just the two known physical
frequencies. Each is bounded using its actual damped denominator in
\eqref{eq:full-fe-lift-fourier}. The remaining terms are bounded on
$|X_i|,|Y_i|\leq0.1+\epsilon_K$ using every fitted quartic coefficient and
the complete trigonometric remainder bounds
$|\sin z-z|\leq|z|^3/6$ and $|1-\cos z|\leq z^2/2$ for real $z$.

For $\beta=\alpha/2$ and
$\omega_{d,j}=\sqrt{\omega_j^2-\beta^2}>0$, the impulse response
$g_j(t)=e^{-\beta t}\sin(\omega_{d,j}t)/\omega_{d,j}$ obeys
\begin{equation}
 \|g_j\|_{L^1(0,\infty)}\leq\frac{2}{\alpha\omega_{d,j}},
 \qquad
 \|g_j'\|_{L^1(0,\infty)}\leq\frac2\alpha+\frac1{\omega_{d,j}}.
 \label{eq:full-fe-green-bounds}
\end{equation}
These estimates include the entire nonlinear forcing remainder and its
Fourier tail. Let $P_j,V_j$ be the resulting all-phase
bounds on the graph's position and velocity. Starting high mode $j$ at
zero adds at most
\begin{equation}
 T_j=P_j+\frac{V_j+\beta P_j}{\omega_{d,j}}
 \label{eq:full-fe-transient}
\end{equation}
to its position envelope. This follows from the exact homogeneous
solution initialized by $(-a_j^*(0),-\dot a_j^*(0))$.
The bounds sum all $62$ high modes and use the cubic Bernstein
control rows of every Hermite element to convert these modal envelopes
into bounds on every centreline point and every phase.

For the two-mode COMSOL integral fit in
Section~\ref{sec:validated-integral-learning}, which satisfies the
nonzero-loss and finite-amplitude KAM conditions, 256-bit Arb
arithmetic encloses $h_{YY}\succeq0.9999999983I$ on the real $0.3$ box.
The complete response graph has centreline amplitude below $1.922696$ mm.
Including the zero-high-mode startup yields an all-time bound below $1.970263$ mm,
strictly inside the $4$ mm net corridor. The latter envelope gives
$|\vartheta|<0.155261$ and $\cos\vartheta>0.987971$, which closes the
feedback-domain bootstrap. The required transverse field is below
$0.921665$ mT and the gradient spectral norm below
$0.02572538\,\mathrm{T\,m^{-1}}$ at every phase and throughout that
startup. These bounds validate feasibility under a $20$ mT transverse
field limit and a $1\,\mathrm{T\,m^{-1}}$ gradient limit; they are
derived bounds, not the extrema of a simulated time history.
Applying the same computation to the fit trained on complete COMSOL FE
trajectories in Section~\ref{sec:comsol-identification} gives graph and
zero-high-state startup bounds of $1.937308$ mm and $2.014078$ mm,
respectively.

The theorem applies to initial low states on the exact certified torus.
The numerical Fourier initial state approximates it within the
certified KAM correction, and finite-time integration tests the feedback
implementation. The theorem uses the exported modal FE coefficients and
fitted weights; independent COMSOL comparisons assess the corresponding
nodal mechanical implementation.
To make the modal definition explicit, if $\Phi$ is the stored full
shape basis, $M,K_b$ are the independently assembled nodal matrices and
$D_\omega=\operatorname{diag}(\omega_j)$, outward Arb products give
\begin{align}
 \|\Phi^{\mathsf T}M\Phi-I\|_\infty&<1.304685\times10^{-14},\\
 \|D_\omega^{-1}\Phi^{\mathsf T}K_b\Phi D_\omega^{-1}-I\|_\infty
 &<3.000518\times10^{-9}.
\end{align}
These are relative kinetic- and stiffness-energy form bounds. They
quantify the relationship between the frozen modal model and the
original assembled FE representation; they are not used to assert
an exact conjugacy or an unproved persistence result for perturbed
matrices. The two-mode shape map used for identification differs from the
reassembled map by at most $1.258\times10^{-14}$ m, and the reference
frequency $\omega_0$ is unchanged.
